%% ============================================================
%% MODELO DE ATIVIDADE · ENSINO E FORMAÇÃO DE MEMBROS
%%
%% Copie este arquivo para atividades/ (atividade-02.tex, por exemplo),
%% inclua-o no relatorio-aacc.tex com \input{atividades/atividade-02}
%% e escreva. Cada \preencher{…} é uma pendência: troque-o pelo seu
%% texto. Cada \figuraprovisoria também: troque-a pela figura de verdade.
%%
%% Uma atividade é UMA entrega sua, com começo, meio e fim. "Participei
%% do projeto X por dois anos" não é uma atividade: são várias. Divida.
%% ============================================================
\begin{atividade}{Treinamento T para os trainees de 2025}{
  tipo         = ensino,
  modalidade   = {},   % a chave da modalidade pedida: EXT, PES, COMP, VPC, EVT, PRD, ORG, ENS, GES, CUR, OUT
  alternativa  = {},   % a modalidade alternativa, se o Colegiado não aceitar a primeira
  horas        = {},   % as horas desta atividade, com base nos registros (Parte IV)
  inicio       = {},   % AAAA-MM
  fim          = {},   % AAAA-MM, ou: em andamento
  projeto      = {},
  papel        = {},   % Responsável, Corresponsável ou Colaborador(a)
  supervisor   = {},   % quem confirma: nome e cargo (o supervisor do projeto, o gerente)
  registros    = {},   % os códigos no SOMA, com ponto e vírgula: NBL-14; NRO-PRO-003-7
  competencias = {},   % as competências das DCN que ela desenvolveu (I a VIII): I, III, V
  disciplinas  = {},   % as disciplinas do seu curso ligadas a ela: ELE075 Eletrônica; ELE067 Circuitos
}

\begin{contexto}
\preencher{O que a equipe precisava aprender, e por quê.}
\end{contexto}

\begin{contribuicao}
\preencher{O que você preparou e ministrou.}
\end{contribuicao}

\begin{desenvolvimento}
\fichatecnica{
  formato       = {},   % Formato (treinamento, tutoria, monitoria)
  treinamentos  = {},   % Treinamentos no SOMA (códigos) — opcional
  publico       = {},   % Público
  cargahoraria  = {},   % Carga horária ministrada
  participantes = {},   % Participantes
  material      = {},   % Material produzido
}

\preencher{O material (módulos, apostila, vídeos, exercícios), a revisão dele e a verificação de conhecimento.}
\end{desenvolvimento}

\begin{resultados}
\preencher{Quantas pessoas, quantas horas, a nota média na verificação, e o que mudou na equipe depois.}
\end{resultados}

\begin{evidencias}
% Uma linha por evidência: \evidencia{tipo}{identificador}{o que mostra}{onde conferir}
% \evidencia{Treinamento}{NRO-TRE-<n> Rev. A}{O treinamento de sua autoria, publicado no SOMA}{SOMA › Treinamentos}
\end{evidencias}

\begin{aprendizados}
\preencher{O que ensinar ensinou a você.}
\end{aprendizados}
\end{atividade}
